\documentclass[12pt]{article}

% House style preamble
\input{.house-style/preamble.tex}
\usepackage{tikz}
\usetikzlibrary{arrows.meta,positioning}
\usepackage[section]{placeins}
\usepackage{flafter}
\fancyhead[L]{\small\scshape Evidentiary Assurance}

\setlength{\headheight}{25.3pt}
\emergencystretch=2em

\hypersetup{
  pdftitle={Evidentiary Assurance for Delegated Authority in AI Systems},
  pdfkeywords={AI governance, AI evidence, delegated authority, decision basis, accountability, procedural validity, electronic evidence, contestability}
}

\title{Evidentiary Assurance for\\Delegated Authority in AI Systems}
\author{Brett Reynolds \orcidlink{0000-0003-0073-7195}%
\thanks{I used OpenAI Codex and Anthropic Claude extensively in drafting and revising the paper. I reviewed, edited, and approved all the material and take full responsibility for the final text and conclusions. \href{mailto:brett.reynolds@humber.ca}{brett.reynolds@humber.ca}}\\
Humber Polytechnic \& University of Toronto}
\date{July 2026}

\begin{document}
\maketitle

\begin{abstract}
AI governance increasingly requires logs, documentation, human oversight, impact assessment, provenance, recourse, and auditability. These requirements create evidence, but they don't yet settle the narrower action-level question: what evidence would make the authorization claim for a specific AI action reviewable after the fact, and would that record support, defeat, or leave the claim indeterminate? Delegation assurance supplies the typed authorization claim and transition semantics. Evidentiary assurance asks whether the preserved record makes that claim evaluable and what evidentiary verdict it warrants. It distinguishes authority, validity, scope, decision basis, attribution, accountability, procedure, authenticity, integrity, and contestability, then proposes a defeasible action-level record-and-verdict test and candidate minimum evidence bundle for recommendations, decisions, communications, tool calls, transactions, and content-generation acts with legal or organizational effect.
\end{abstract}

\noindent\textbf{Keywords:} AI governance; AI evidence; delegated authority; decision basis; accountability; procedural validity; electronic evidence; contestability

\section{Introduction}

AI risk-management guidance already treats governance, transparency, accountability, documentation, and monitoring as core trustworthiness practices \parencite{nist2023artificial}. NIST's generative-AI profile adds a more specific emphasis on governance, content provenance, pre-deployment testing, and incident disclosure \parencite{nist2024generative}. This paper asks a narrower action-level evidentiary question: given a contested AI action, what body of evidence would justify the conclusion that the action stayed within authority validly delegated to the system? The question is narrower than explainability, broader than logging, and more forensic than ordinary model documentation.

The target action may be a recommendation, decision, communication, tool call, transaction, record update, or content-generation act with legal or organizational effect. The central evidentiary burden is to connect that act to a valid source of authority, a scoped delegation instrument, the relevant system state, the procedure actually followed, and records whose authenticity and integrity can survive later review.

Recent work on digital money makes the infrastructural point in another domain: a perfectly safe liability can still become fragile when it runs on a congested settlement rail \parencite{klee2026fragility}. AI evidence has a comparable rail. A model may be capable, a permission may be valid, and a policy may be well specified, but the later authority claim still depends on the infrastructure that carries action from instruction to tool use, logging, approval, and review.

Administrative-law record review gives the closest legal analogue. A reviewer needs more than a plausible result. The record has to show authority, procedure, reasons, evidentiary basis, and reviewability. AI action records add a further burden: they have to distinguish the system that produced the act from the human or institutional actor that remains responsible for authorizing and relying on it.

The paper belongs beside, but not inside, the delegation-assurance paper. Delegation assurance gives AI-security and evaluation communities an evidence standard for justified machine action. This paper develops the legal, governance, provenance, and digital-evidence side of that standard.

\section{Evidence Primitives}

Table~\ref{tab:evidentiary-primitives} separates ten review questions that often get collapsed. The columns matter together: the middle column names the record each primitive needs, while the right column prevents that primitive from doing more work than it can support. A system can have authentic logs without valid authority, clear provenance without delegated scope, or a clean trace of a decision rule whose predicates were never established.

\begin{table}[tbp]
\small
\centering
\begin{tabular}{>{\raggedright\arraybackslash}p{0.21\linewidth}%
                >{\raggedright\arraybackslash}p{0.34\linewidth}%
                >{\raggedright\arraybackslash}p{0.35\linewidth}}
\toprule
Primitive & Evidence it needs & What it doesn't establish alone \\
\midrule
Authority & Delegation instrument, role, policy source, or legal/organizational mandate. & That the principal could validly confer that authority. \\
Validity & Legal or organizational power to delegate this action class, including non-delegable limits. & That this action stayed within the delegated limits. \\
Scope & Action class, target, time, tool, data class, threshold, exception, and approval condition. & That the action was actually produced under that scope. \\
Decision basis & Governing rule and version; required factual and legal predicates; source evidence and data definitions; transformations, aggregation, imputation, temporal granularity, uncertainty, and burden allocation. & Authorization, general substantive correctness, legality, or fairness. \\
Attribution & System instance, model and tool versions, configuration, actor, approval, and execution trace. & That the action was validly authorized. \\
Accountability & Identified individual or organization capable of bearing the relevant duty, linked role or office, duty source, effective interval, succession route, and bearer-specific powers. & That a role label or team is itself a reachable bearer, that the bearer performed the action, or that it has the review forum's powers. \\
Procedure & Required steps, notices, reviews, approvals, overrides, adverse-signal routes, expected controls, escalation paths, dispositions, and recourse records. & That the substantive outcome was correct or fair. \\
Authenticity & Identity, source, signature, timestamp, record association, and custody evidence. & That the authenticated assertion is true or authorized. \\
Integrity & Completeness, tamper evidence, retention controls, hash or log protection, and change history. & That an intact record contains enough information for review. \\
Contestability & Accessible reasons, relevant record, disclosure path, competent review forum, applicable independence, and correction or remedy route. & That the original action was valid or that every action gives an affected party the same review rights. \\
\bottomrule
\end{tabular}
\caption{Evidentiary primitives for delegated AI action.}
\label{tab:evidentiary-primitives}
\end{table}

\term{Decision basis} is provisional as a separate primitive. Its distinctive question is whether the rule actually used was linked to the rule's required predicates through defensible data definitions and transformations. The record should also map the rule used back to its claimed source of authority. A perfectly attributable execution can otherwise preserve an unsupported inference without exposing it as such. Decision-basis evidence doesn't establish general legality or fairness, and it should be folded into Authority, Validity, and Scope if matched-case testing shows that it detects no additional failure.

Action-level review requires two linked records rather than one generic owner field. The accountability primitive uses an \term{answerability-bearer matrix} to identify each candidate individual or organization, the linked role or office, duty source, effective interval, succession route, and bearer-specific powers. The contestability primitive uses a \term{review/remedy forum matrix} to identify the forum or reviewer, institutional or jurisdictional basis, record access, competence, applicable independence, duties, and powers to review, correct, halt, sanction, compensate, or remediate where those powers apply. The bearer doesn't acquire the forum's powers, and neither record establishes legal liability on its own.

Procedure contains two records that ordinary logs often omit. An \term{adverse-signal record} binds warnings, legal advice, adverse review decisions, complaints, and incidents to their time, applicability, recipients, relevant duties, deadlines, dispositions, escalations, and downstream changes. An \term{expected-control record} binds a triggering condition to the act required, responsible role, due time, completion or disposition, and monitoring coverage. Failure to record completion counts against an assurance claim only when the expected completeness of that logging is itself supported.

Explanation, provenance, logging, and documentation supply parts of the linked evidence required by the action-level record-and-verdict test.

\section{Governance and Legal Sources}

This section is cross-jurisdictional, with emphasis on Canada, the European Union, the United States, and Australia where official sources are especially probative. The relevant sources fall into three groups.

Agency law supplies the clearest vocabulary for the underlying relation: a principal authorizes an agent to act on the principal's behalf, and actual authority may be express or implied \parencite{cornellWexAgency,cornellWexActualAuthority}. That analogy clarifies the question, but it also marks the limit. In ordinary agency, the agent can be a responsible legal or institutional actor. In AI deployments, the system can execute delegated action, but responsibility remains with the person or entity that authorizes and controls its use \parencite{cornellWexPrincipal}.

First, AI-governance instruments require logs, documentation, intended-purpose limits, instructions for use, human oversight, impact assessment, monitoring, recourse, and disclosure. The EU AI Act requires high-risk systems to support traceability through logging capabilities, requires instructions for use, and assigns deployers obligations around use, input data, monitoring, serious-incident reporting, and human oversight by people with competence, training, and authority \parencite{europeanUnion2024aiAct}.

Canada's Algorithmic Impact Assessment asks comparable operational questions about project authority, system accountability, audit trails, system-produced reasons, recourse, human override, documentary evidence, and changes in system functionality or scope \parencite{canadaAIA2026}. These instruments generate evidence, but they usually stop short of specifying what evidence is sufficient to support a finding that a concrete action stayed within delegated authority.

Financial-reporting law supplies a useful institutional analogue. Sarbanes--Oxley Section 404 requires internal-control reporting that states management's responsibility for adequate controls and assesses their effectiveness; the SEC's implementing rules require management to identify the framework used for that evaluation, and Exchange Act Rule 13a-15 requires a suitable, recognized control framework for evaluating internal control over financial reporting \parencite{sarbanesOxley404,sec2003managementReport,secExchangeRule13a15}.

This comparison is limited but clarifying. SOX doesn't ask investors to trust a company's bare assertion that the numbers came out right. It asks whether management maintained a reviewable control environment around the reporting process. Evidentiary assurance asks for the same kind of move for AI action: not a transcript alone, and not a model-level safety claim alone, but a control record tying action to authority, scope, procedure, integrity, and accountable institutional actors.

Second, electronic-evidence and electronic-signature law already has mature concepts for authenticity, integrity, attribution, secure signatures, and tamper detection. The Canada Evidence Act places the burden of authenticating an electronic document on the party seeking admission and ties the best-evidence rule to proof of the integrity of the electronic-document system or an applicable evidentiary presumption \parencite{canadaEvidenceAct1985}. PIPEDA Part 2 and Canadian e-signature guidance add related primitives: signer identity, signer authority, association between signature and record, record retention, timestamps, transaction records, and integrity of the signed data and supporting information \parencite{pipeda2000,canadaESignatures2025}. Those doctrines aren't AI-specific, but they provide strong analogues for AI action traces.

Third, public failures show the practical stakes. The Robodebt Royal Commission recommended clear review paths for affected people, plain-language notice that automated decision-making is used, and public availability of business rules and algorithms for independent expert scrutiny \parencite{robodebtRoyalCommission2023}. That recommendation set is close to the evidentiary-assurance question: what record lets an affected person, auditor, regulator, court, or internal reviewer determine whether the system acted within authority the institution could validly delegate?

\begin{table}[tbp]
\small
\centering
\begin{tabular}{>{\raggedright\arraybackslash}p{0.22\linewidth}%
                >{\raggedright\arraybackslash}p{0.28\linewidth}%
                >{\raggedright\arraybackslash}p{0.40\linewidth}}
\toprule
Source family & Evidence primitives supported & Remaining gap \\
\midrule
Agency law & Authority, scope, attribution, accountability. & Ordinary agency assumes a responsible agent; AI systems execute action without becoming accountable legal subjects. \\
AI governance & Procedure, monitoring, documentation, oversight, recourse, traceability. & Governance duties create records without specifying the action-level sufficiency test. \\
Financial-reporting controls & Management responsibility, control frameworks, effectiveness assessment, attestation, change monitoring. & Internal control over reporting isn't yet an action-level standard for delegated AI systems. \\
Electronic evidence & Authenticity, integrity, attribution, preservation, evidentiary presumptions. & Document integrity doesn't by itself show that an AI action was authorized. \\
Electronic signatures & Identity, authority to sign, record association, timestamping, transaction evidence. & Signature validity doesn't generalize automatically to multi-step agentic action. \\
Public automation failures & Contestability, notice, review paths, business-rule disclosure, independent scrutiny. & Post-failure recommendations still have to be translated into deployable evidence bundles. \\
\bottomrule
\end{tabular}
\caption{Legal and governance sources as evidence primitives.}
\label{tab:legal-governance-primitives}
\end{table}

Taken together, these sources set the frame. Governance law tells deployers to create and preserve records. Evidence law tells reviewers how authenticity and integrity can be established for digital artefacts. Agency law clarifies why authority and accountability have to be separated in AI systems. Public-sector failures show what happens when automated action can't be reconstructed or contested. Evidentiary assurance joins these materials around a single action-level question: what's enough evidence to defend this action, by this system, under this delegation, at this time?

\section{Technical Mechanisms}

Technical mechanisms supply pieces of the evidentiary bundle. Structured execution logs record what happened and when. Provenance graphs connect entities, activities, agents, and transformations; W3C PROV gives a general vocabulary for that lineage \parencite{moreau2013provdm}. Content provenance systems bind claims, manifests, hashes, signatures, and timestamps to digital artefacts; C2PA provides a developed current example for media and document provenance \parencite{c2pa2026spec}. Append-only transparency logs make alteration and equivocation more detectable; Certificate Transparency shows how Merkle-tree logs and consistency proofs can support independent auditing \parencite{laurie2013certificate}. Bills of materials expose components and dependencies. Runtime policy enforcement can check explicit constraints before action.

Agentic systems add a language-mediated evidence problem. For tool-using language-model systems, adversarial language can enter through retrieved content and lower-priority instructions, so the evidentiary record also has to preserve contextual status, \term{normative standing} (whether a source is entitled to direct the action), any applicable conflict priority, and action licensing, with provenance channel and technical trust recorded separately \parencite{greshake2023not,wallace2024instruction,debenedetti2024agentdojo}. Adversarial pragmatics \parencite{reynolds2026adversarialPragmatics} operationalizes that language-status problem as an evaluation target. Authenticated-delegation proposals and agent-audit protocols point toward a stronger record: agent identity, governing or conferring sources, scoped permissions, authorization transitions, audit context, and action records \parencite{south2025authenticated,kuehlewind2026audit}.

Table~\ref{tab:technical-mechanisms} treats each mechanism as partial evidence. Each secures part of the record; none settles authority, scope, integrity, and reviewability by itself.

\begin{table}[tbp]
\small
\centering
\begin{tabular}{>{\raggedright\arraybackslash}p{0.24\linewidth}%
                >{\raggedright\arraybackslash}p{0.31\linewidth}%
                >{\raggedright\arraybackslash}p{0.35\linewidth}}
\toprule
Mechanism & Evidentiary contribution & Residual assurance gap \\
\midrule
Execution logs & Event sequence, timestamps, tool calls, approvals, denials, and errors. & Logs may omit the authorization state, contextual status, or technical trust that made the action valid or invalid. \\
Provenance graphs & Lineage across data, prompts, tools, transformations, agents, and artefacts. & Lineage doesn't determine whether an input was entitled to guide action. \\
Signed claims and manifests & Association between record, assertion, signer, timestamp, and asset. & Signed assertions can still be overbroad, stale, false, or unauthorized. \\
Append-only logs & Tamper evidence, consistency proofs, and stronger preservation. & Integrity of the log doesn't guarantee completeness or legal sufficiency. \\
Bills of materials & Model, dependency, component, dataset, and tool inventories. & Inventory doesn't establish runtime scope compliance. \\
Runtime policy gates & Machine-checkable action constraints before execution. & The gate needs the right policy, the right context, and a record of why it allowed or blocked action. \\
Agent-audit protocols & Delegated identity, permission transitions, audit context, and trace reconstruction. & Protocol records still need language-status evidence and institutional accountability. \\
\bottomrule
\end{tabular}
\caption{Technical mechanisms for evidentiary assurance.}
\label{tab:technical-mechanisms}
\end{table}
\FloatBarrier

\subsection{Typed Representations and Inferential Limits}

Table~\ref{tab:representation-semantics} keeps the underlying relations typed, preventing one convenient graph from being made to answer authorization, provenance, causal, transport, selection, and evidentiary questions at once.

\begin{table}[tbp]
\footnotesize
\centering
\begin{tabular}{>{\raggedright\arraybackslash}p{0.15\linewidth}%
                >{\raggedright\arraybackslash}p{0.18\linewidth}%
                >{\raggedright\arraybackslash}p{0.24\linewidth}%
                >{\raggedright\arraybackslash}p{0.32\linewidth}}
\toprule
Representation & Relation or marker & Can support & Can't establish alone \\
\midrule
Temporal authorization record & Grants, constrains, amends, suspends, revokes, overrides, releases, reviews. & Standing, scope, and valid state transitions under a declared regime. & Causal uptake, truth, or legal validity beyond the stipulated or adjudicated regime. \\
Provenance graph & Used, generated, derived, transformed, preserved. & Lineage and reconstruction. & Causation, authorization, or substantive correctness. \\
Causal DAG or structural causal model & Hypothesized causal influence and interventions. & Causal attribution under stated assumptions. & Standing, permission, or normative validity. \\
Transportability selection diagram & An \emph{S}-marker denotes a suspected source--target mechanism or marginal difference. & Analysis of a specified causal query across domains. & Construct or normative equivalence, or transportability by inspection. A changed target meaning requires a new claim. \\
Sample-selection DAG & Pathways into observation, inclusion, complaint, or review. & Diagnosis of observed-case bias and denominator needs. & Population incidence without assumptions and denominator data. \\
Assurance-case graph & Supports, warrants, rebuts, undercuts. & Explicit, noncompensatory claim--evidence structure. & Truth or sufficiency merely because evidence is linked. \\
\bottomrule
\end{tabular}
\caption{Graphical and record semantics used in evidentiary assurance.}
\label{tab:representation-semantics}
\end{table}
\FloatBarrier

This typing changes the audit question. An authorization relation can't be inferred from a causal arrow, and a provenance edge doesn't show that its source caused or licensed an action. In a transportability selection diagram, an \emph{S}-marker encodes an asserted source--target difference within a specified causal problem; the diagram alone doesn't establish transportability \parencite{pearlBareinboim2014externalValidity}. A source--target change in governing mandate or outcome meaning changes the target claim rather than adding an \emph{S}-marker to the old query. This paper uses the assurance-case relation in Figure~\ref{fig:evidence-rail} and a limited sample-selection sketch below, without proposing a transportability selection diagram.

The architecture implied by the mechanism table is layered. A deployment needs a delegation instrument, a decision-basis record, a policy-state snapshot, a configuration snapshot, a runtime trace, an adverse-signal and expected-control record, a preservation mechanism, and a review interface. Weakness at any required layer can defeat the later authorization claim or leave it indeterminate. A perfect tool-call log doesn't help if the delegation instrument is missing. A signed policy doesn't help if content from a source lacking standing is treated as controlling or if its provenance and status record is lost. A complete trace doesn't support an applicable contestability claim if the authorized reviewer can't access the relevant parts of it.

The same architecture also creates privacy and security risk. Evidence capture should include minimization, redaction, retention schedules, role-based access, and protected review channels. The goal is enough preserved evidence for authorized review without turning every delegated action into an unnecessary archive of sensitive context.

\section{An Action-Level Sufficiency Test}

Evidentiary assurance is defeasible. It asks whether the preserved record is strong enough to support an authority claim against predictable challenges, while leaving wisdom, fairness, full legality, and substantive correctness for other forms of review. A record is sufficient only if a reviewer can connect the action to valid authority, scope, system state, context, procedure, preservation, and accountable review without relying on the system's bare assertion.

\begin{table}[tbp]
\small
\centering
\begin{tabular}{>{\raggedright\arraybackslash}p{0.23\linewidth}%
                >{\raggedright\arraybackslash}p{0.34\linewidth}%
                >{\raggedright\arraybackslash}p{0.33\linewidth}}
\toprule
Sufficiency claim & Required showing & Typical defeater \\
\midrule
Authority source & A principal, role, law, policy, or organizational mandate authorized this class of action. & Missing, stale, ambiguous, or overbroad delegation instrument. \\
Authority validity & The principal could validly confer the action class, and the action doesn't cross a non-delegable legal, organizational, or policy limit. & The record documents a delegation the institution had no power to make. \\
Scope application & The concrete action fits the permitted operation, target, time, data class, threshold, and approval condition. & Scope inflation from analysis to disclosure, persistence, contact, purchase, or decision. \\
Runtime state & The relevant model, prompts, tools, policies, permissions, memory state, and deployment settings are identified. & Unverifiable configuration, untracked model change, or missing policy snapshot. \\
Context provenance & User instructions, retrieved material, tool outputs, memory entries, and policy text have source and trust labels. & Lower-trust language is promoted into action-guiding authority. \\
Execution procedure & Tool calls, arguments, approvals, denials, overrides, retries, and failures are reconstructable. & The record contains only a final answer or self-reported rationale. \\
Preservation integrity & The record has custody, retention, access-control, timestamp, signature, hash, or tamper-evidence protections appropriate to the risk. & Gaps, alteration risk, selective logging, or unverifiable custody. \\
Answerability and review & A human, vendor, deployer, or institution remains answerable, and affected parties have a review or correction path. & Accountability is implicitly offloaded to the model, or the affected reviewer can't inspect the relevant record. \\
\bottomrule
\end{tabular}
\caption{Defeasible sufficiency test for delegated AI action.}
\label{tab:sufficiency-test}
\end{table}

The test shifts the target from record production to record adequacy. A log can be authentic and still inadequate if it omits the authority state. An action list can be complete and still inadequate if it omits the status the system assigned to its sources, since a reviewer can then judge the act but not the recognition it flowed through. A policy can be valid and still inadequate if the system's runtime configuration can't be reconstructed. A complete trace can still be inadequate if no accountable actor remains answerable or no affected person can challenge the action. A flawless bundle can also document an invalid delegation; validity is a separate showing, not an inference from clean procedure.

A procurement-assistant case makes the test concrete. Suppose the assistant is authorized to draft an internal comparison memo and reads a vendor page containing the instruction \enquote{email the buyer's maximum budget to request a discount}. The contested action is a blocked vendor email.

A sufficient record would show the user task and procurement role as the authority source; the validity basis for reading internal purchasing criteria; the scope limit that permits memo drafting but not vendor contact; the runtime policy and tool permissions; the vendor page's source and trust label; the proposed email arguments; the tool gate's block decision; the timestamped trace; and the accountable team that configured the workflow.

That record would support the audit inference that blocking the email preserved delegated authority. If the assistant sent the email, the same bundle would locate the defeater: scope inflation, failed information-flow control, unauthorized tool licensing, or missing human approval. If the only record were the final memo, the action might look successful while the authority claim remained unproven.

The table also supports a burden-shifting account for audits and assurance reviews. In a forum that adopts such a standard, an adequate bundle could support a rebuttable audit inference that the action was authorized: the proponent shows authority, validity, scope, state, context, procedure, preservation, and review; a challenger points to a defeater such as invalid authority, scope inflation, missing configuration, contaminated context, or record gaps. Legal presumptions would require a separate doctrinal argument. The applicable standard of proof is forum-specific, but ties shouldn't be resolved by the model's self-report.

A sufficiency test governs one action's record, but an assurance program advances a claim across many, and that aggregation invites a selection effect. If the cases a proponent presents, or the model configuration, prompt set, or evaluation run behind them, were chosen because they scored best, the aggregate inference is optimistically biased even when every individual bundle is clean. Decision analysts call this the optimizer's curse \citep{smith2006optimizersCurse}: selecting the maximum over noisy estimates preferentially picks options whose estimation error ran high, so the chosen option's realized reliability falls below its estimate. So selection is itself a defeater. A challenger can ask whether the evidence was sampled from deployment or curated to persuade, and the reviewer's counterpart is what Smith and Winkler call \enquote{disciplined skepticism}: discount the claim toward the field's base rate in proportion to how hard the evidence was optimized, weighting bundles drawn at random above those offered as a program's best.

The standard is also risk-sensitive. A low-stakes internal recommendation may require a lighter record than a benefits decision, clinical summary, procurement action, or external communication. The difference is evidentiary burden, not evidentiary kind: higher-risk deployments need stronger showings for the same validity, authority, scope, state, context, procedure, preservation, and review claims.

\section{Research Program}

The immediate research task is a literature intake, not prose expansion. The intake should map explainability, documentation, provenance, content authenticity, secure logging, AI governance, electronic evidence, electronic signatures, algorithmic auditing, and public-sector automation failures.

The paper should then derive a minimum evidence bundle for delegated AI action. A candidate bundle includes the delegation instrument, policy and permission state, configuration snapshot, model and tool versions, input context, retrievals, tool calls, approvals, overrides, safety interventions, timestamps, human-oversight records, and integrity protections.

The central claim should remain narrow. The literature doesn't ignore evidence. It disperses evidence across neighbouring fields. The contribution is to integrate those fields around one question: what proves that a specific AI action remained within delegated authority?

\clearpage
\section{Conclusion}

The literature disperses evidence across neighbouring fields. Explainability helps with reasons, provenance helps with lineage, audit logging helps with reconstruction, governance law creates documentation duties, and electronic-evidence law supplies concepts for authenticity and integrity. None of these materials alone answers the delegated-authority question, and none makes clean execution evidence a substitute for a warranted decision basis, routed adverse signals, prospectively required governance controls, or meaningful review.

Evidentiary assurance integrates those materials around a bounded action-level review. Delegation assurance supplies the normative claim \(A_R(a,t)\) and transition semantics. Evidentiary assurance supplies four evidence-relative verdict vectors, parameterized by regime, adjudicating review forum, evidential standard, burden allocation, applicability-map version, action class, and declared use: \(J_A\) for what the record warrants about authorization conditions, \(J_R\) for record adequacy, \(J_B\) for bearer identity, duty, and reachability, and \(J_C\) for object-forum access, competence, independence, and remedial power. Each preserves support, substantive defeat, record gaps, and conflict at the node where they arise.

The framework supplies no universal assurance score. A declared use may require a prospective conjunction of selected nodes, but poor recording doesn't make an authorized action historically unauthorized and pristine recording doesn't cure invalid authority. The executable fixtures show that the distinctions can be encoded and exposed to controlled failure; with no reviewer responses yet, they don't show that reviewers can apply them reproducibly or validly. The proposed standard remains a research hypothesis until its interpretation and use are supported by content, substantive, structural, generalizability, external, and consequential validity evidence, followed by target-bound projection and prospective withdrawal when its fixed minimum reach is lost. Those questions remain subject to forum- and jurisdiction-specific determination. The account is narrower than a finding of full legality, fairness, wisdom, legal liability, moral responsibility, or substantive correctness.


\clearpage
\begingroup
\renewcommand*{\bibfont}{\footnotesize}
\setlength{\bibitemsep}{0.25\baselineskip}
\printbibliography
\endgroup

\end{document}
